\section{Multi-Domain Generalization}

\label{sec:domains}

A key architectural claim of the UAH is that the seven-layer harness is
\textbf{domain-agnostic}: the same infrastructure serves agents across
radically different task types, with domain specificity concentrated in the
Layer 3 agent specifications (skills, tools, instructions) rather than in
the infrastructure layers.

We formalize this claim:

\begin{definition}[Domain-Agnostic Harness]
A harness $H$ is \emph{domain-agnostic} if for any two task domains
$\mathcal{D}_1$ and $\mathcal{D}_2$, the infrastructure layers $L_1, L_2,
L_4, L_5, L_6, L_7$ are identical and only $L_3$ (Agent Runtime)
configurations differ.
\end{definition}

We demonstrate this property across four deployment domains.

\subsection{Software Engineering Domain}

The software engineering agent (``Hanzo Dev'') uses:
\begin{itemize}[leftmargin=1.5em]
    \item \textbf{L3 configuration:} Orchestrator agent + specialist agents
          for planning, code generation, testing, and review.
    \item \textbf{Primary tools:} Code intelligence (10 tools), shell
          execution (8 tools), filesystem (12 tools).
    \item \textbf{Default tier:} $T_1$ (Terminal Pod). Upgrades to $T_3$ for
          CI/CD tasks requiring cloud services.
    \item \textbf{Improvement signal:} Test pass rate, linting score, PR merge
          acceptance rate.
\end{itemize}

Performance on SWE-bench Verified: 74.3\% resolved rate.

\subsection{Marketing Automation Domain}

The marketing agent (``Hanzo Growth'') uses:
\begin{itemize}[leftmargin=1.5em]
    \item \textbf{L3 configuration:} Research agent + copywriter agent +
          analytics agent + publisher agent.
    \item \textbf{Primary tools:} Web search, browser automation, data
          analysis, email/CMS publishing tools (via custom MCP extensions).
    \item \textbf{Default tier:} $T_0$ (Shared Gateway) for research and
          copywriting. $T_2$ (Operative Desktop) for browser-based campaign
          management.
    \item \textbf{Improvement signal:} Click-through rate, conversion rate,
          campaign ROI.
\end{itemize}

The marketing domain illustrates the Channel Integration layer's value:
campaigns are triggered from Slack, executed autonomously, and results are
reported back to Slack and the connected Linear project.

\subsection{DevOps Operations Domain}

The DevOps agent (``Hanzo Ops'') uses:
\begin{itemize}[leftmargin=1.5em]
    \item \textbf{L3 configuration:} Alert triage agent + diagnosis agent +
          remediation agent, connected to PagerDuty webhook via Layer 4.
    \item \textbf{Primary tools:} kubectl, Helm, Prometheus queries, log
          analysis, GitHub (deploy workflows).
    \item \textbf{Default tier:} $T_3$ (Cloud VM with direct cluster access).
    \item \textbf{Improvement signal:} Mean time to resolution (MTTR),
          false positive alert rate, successful remediation rate.
\end{itemize}

Over 90 days of production deployment, Hanzo Ops reduced mean MTTR from 47
minutes to 8.3 minutes for the class of incidents within its capability scope
(database performance issues, deployment failures, certificate expirations,
and capacity alerts).

\subsection{Research Domain}

The research agent (``Hanzo Research'') uses:
\begin{itemize}[leftmargin=1.5em]
    \item \textbf{L3 configuration:} Literature search agent + experimental
          design agent + data analysis agent + writing agent.
    \item \textbf{Primary tools:} ArXiv/PubMed search, PDF parsing, statistical
          analysis tools, Python/R execution, citation management.
    \item \textbf{Default tier:} $T_3$ (Cloud VM for compute-heavy analysis).
          Upgrades to $T_4$ (Local Desktop) when GPU acceleration is required.
    \item \textbf{Improvement signal:} Literature coverage completeness,
          statistical validity (verified by Reviewer agent), citation accuracy.
\end{itemize}

\begin{table}[H]
\centering
\caption{Domain-Specific Performance Metrics}
\label{tab:domains}
\begin{tabular}{lllcc}
\toprule
\textbf{Domain} & \textbf{Primary Metric} & \textbf{UAH Score} &
\textbf{Default Tier} & \textbf{Improvement Loops} \\
\midrule
Software Eng.    & SWE-bench resolved rate & 74.3\% & $T_1$ & 4 \\
Marketing Auto.  & Campaign setup accuracy & 81.2\% & $T_0$/$T_2$ & 3 \\
DevOps Ops.      & Incident resolution rate & 76.8\% & $T_3$ & 4 \\
Research         & Literature coverage      & 88.4\% & $T_3$/$T_4$ & 4 \\
\midrule
\multicolumn{2}{l}{\textbf{Mean}} & \textbf{80.2\%} & --- & --- \\
\bottomrule
\end{tabular}
\end{table}

Across all four domains, the infrastructure layers (L1, L2, L4, L5, L6, L7)
are identical. Total additional code required to deploy a new domain is
approximately 400--800 lines: the Layer 3 agent specification, any custom
MCP tool extensions, and the domain-specific improvement signal definition.

% ============================================================================
% 8. PRODUCTION IMPLEMENTATION
% ============================================================================
